\section{从已暴露失败到第二轮研究}
\label{sec:revision-design}

第一轮提供了一个真实但尚不充分的研究闭环。它发现的三个缺口直接决定第二轮问题：稀疏观测的三折训练集不足六条，导致运动模型被结构性排除；静态均值点估计的10\%容忍条件只在开发验证中满足，未在首次测试保持；总体边际区域掩盖运动异常子群的较低覆盖。第二轮因此研究折数、选择门槛及稳健组件的组合，并比较两种预先声明的区域校准方式。修订不以消除所有失败作为交付条件，最终结论服从新数据。

\subsection{旧结果是先验，新测试才是验证}
第二轮把第一轮结果作为已经暴露的先验，使用独立运行目录\path{runs/revision_reference}。发现集和验证集均包含四个常规场景及稀疏场景，种子分别为1646191与1657383，每场景分别20例与30例。第一轮数值源码、轨迹与冻结记录保持原字节；修订工具通过子类和独立入口扩展，实现与历史证据分开核验。

第二轮的全新主测试包含720例，压力测试400例，观测数边界测试420例，独立校准800例，覆盖检查1200例。机器可读方案\path{research/revision_plan.json}及评价器源码在新研究初始化时一并记录摘要。新案例只在方法冻结、源码和轨迹核验通过之后生成。因此，第二轮不能被描述成对第一轮研究路线的独立重复；它是\emph{依据已知失败作修订，再用新数据检验}的后续研究。

\subsection{更严格的工具约束}
第二轮显式开放原求解器已有的\texttt{cv\_folds}字段，让模型能够把折数写入可复核的方法定义。候选、单组件消融、诊断后修订、独立数学检查和验证证据分别留痕；冻结前要求候选与第一轮冻结方法及预设基线具有验证记录。工具也修订了第一轮消融说明继承父文本的问题：新登记同时保留父定义、实际改动和最终执行配置，避免把父算法的数学文字当成子算法的准确说明。

这些约束增加了过程可审查性，但并不证明规划器优于其他研究策略。两轮研究目标、开发数据和搜索空间不同；两次轨迹不能合并成同条件的两次随机重复，也不构成模型间的等预算对照。

\subsection{观测数与折数的结构关系}
若$K$个时间交错验证折的大小相差至多一条，最小训练折有
\[
 n_{\mathrm{train,min}}=n-\left\lceil\frac{n}{K}\right\rceil
\]
条观测。原实现要求运动训练折至少六条，因此必要计数条件为$n-\lceil n/K\rceil\ge6$。$n=8$时，$K=3$产生五条训练观测，而$K=4$产生六条。该变化只消除了计数障碍；时间跨度、几何秩、求解边界和噪声仍决定模型能否可靠拟合。

增加折数还会缩小每个验证折并增加拟合次数。实际评分汇总全部留出观测，故不能仅凭单折变小就断言总体评分更受异常影响；折数同时改变训练样本和预测误差的依赖关系，其净影响需要验证。更多拟合也提高计算成本。为避免仅展示有利的$n=8$，预登记边界测试完整覆盖$n=6,7,8,9,10,11,12$，在静态与运动无污染场景各生成30例，并保留所有不利结果。

\subsection{按可观测选模结果分组校准}
第一轮的子群欠覆盖提示误差分布异质性。本轮使用已有的条件共形框架\cite{vovk2012conditional}，不将分组校准称为原创方法。记完整观测为$Z$，冻结估计器为$\widehat x(Z)$，正尺度为$g(Z)$，并定义可观测分组
\[
 G(Z)\in\{\text{静态输出},\text{恒速输出}\}.
\]
它由冻结求解器的实际选择给出，不读取真实运动类别、污染标签或当前位置。对组$j$，取校准分数$S_i=\|X_i-\widehat x(Z_i)\|_2/g(Z_i)$中满足$G(Z_i)=j$的$m_j$项，并令
\[
 k_j=\left\lceil(m_j+1)(1-\alpha)\right\rceil,\qquad
 q_j=\begin{cases}S^{(j)}_{(k_j)},&k_j\le m_j,\\+\infty,&k_j>m_j.\end{cases}
\]
新区域为$C(Z)=\{x:\|x-\widehat x(Z)\|_2\le q_{G(Z)}g(Z)\}$。空组也按无界区域处理，不能用零半径或\texttt{null}表示已经获得有效的有限区域。

在估计器、尺度和分组函数预先冻结，校准与新案例独立同分布的条件下，给定新案例所属组及校准组大小，组内分数的交换性给出秩覆盖界。无并列时，新分数在$m_j+1$项中的秩均匀；有并列时采用不严格包含保持保守性。因此
\[
 \Pr\{X_{\rm new}\in C(Z_{\rm new})\mid G(Z_{\rm new})=j,\,M_j=m_j\}\ge1-\alpha.
\]
这个概率仍对校准分数取平均，不能解释为固定本次校准集后的保证。实际输出的静态组也不是“真实静止组”：误选模型会把不同物理场景混在一起。压力分布、边界分布、按半径拒绝的子集和任意精细真值场景均不自动获得这一条件覆盖。

本轮在相同800例校准与1200例覆盖数据上，同时报告全局与分组方法的覆盖、半径、接受率及真实场景内结果。分组样本减少可能增大分位数不确定性，改善覆盖也可能以更大半径为代价；选择最漂亮的单一指标不足以评价区域是否有用。
